\cleardoublepage

\chapter*{Resumen\markboth{Resumen}{Resumen}}

El sector TIC genera entre el 2 y el 4\,\% de las emisiones globales de CO$_2$, pero la formación profesional en informática apenas incorpora la eficiencia energética como criterio de calidad. Este TFM diseña una propuesta de integración transversal de Green IT: cinco prácticas de laboratorio (P0--P4) en el módulo MP0485 Programación del CFGS DAM (1.º curso), sin alterar el currículo oficial. La secuencia lleva al alumnado desde la concienciación sobre la huella digital hasta la medición, optimización y toma de decisiones carbon-aware sobre su propio código, con CodeCarbon y la API Electricity Maps, sustentada en el aprendizaje experiencial de Kolb y el aprendizaje basado en proyectos.

\vspace{0.5cm}
\noindent\textbf{Palabras clave:} Green IT, sostenibilidad digital, educación para el desarrollo sostenible, Formación Profesional, programación, huella de carbono del software, CodeCarbon.

\chapter*{Abstract\markboth{Abstract}{Abstract}}

The ICT sector accounts for 2--4\,\% of global CO$_2$ emissions, yet vocational computing education rarely treats energy efficiency as a software quality criterion. This thesis designs a cross-curricular Green IT proposal: five laboratory practices (P0--P4) embedded in module MP0485 Programming of the Higher VET programme in Multiplatform Application Development (first year), without altering the official curriculum. The sequence guides students from digital-pollution awareness to measurement, optimisation and carbon-aware decision making on their own code, using CodeCarbon and the Electricity Maps API, grounded in Kolb's experiential learning cycle and project-based learning.

\vspace{0.5cm}
\noindent\textbf{Keywords:} Green IT, digital sustainability, education for sustainable development, vocational education, programming, software carbon footprint, CodeCarbon.
